
\section{HERMES neural network model for amino acid propensities}
 HERMES is a self-supervised structure-based machine learning model for
 proteins~\cite{Visani2024-sh,Pun2024-qz}. HERMES has a 3D rotationally equivariant architecture and
 is trained to predict a residue's amino acid identity from its surrounding atomic environment
 within a 10 \AA~ of the central residue's C-$\alpha$ in the 3D structure. The input to the neural
 network is a point cloud of atoms from a given structural neighborhood, in which all atoms
 associated with the central residue are masked. These point clouds are featurized by atom types,
 including computationally-added hydrogens, partial charge, and Solvent Accessible Surface Area
 (SASA). The point clouds are first projected onto the orthonormal Zernike Fourier basis, centered
 at the position of the (masked) central residue's C-$\alpha$ (Fig.~1A). The resulting
 \textit{holographic} projections are fed to a stack of SO(3)-equivariant layers, the output of
 which is SO(3)-invariant, and is then passed through a multi-layer perceptron (MLP) to generate the
 desired predictions. The MLP's outputs are termed pseudo-energies, associated with the 20 different
 possible amino acids that can be placed at the masked residue, and act as logits within a softmax
 function to compute probabilities of each amino acid.

Each HERMES model is an ensemble of 10 individually-trained architectures. In this work, we used
HERMES models trained on structural data with different levels of added noise: 0.00 noise (the
original PDB structure), and 0.50 noise, which is obtained by adding Gaussian noise to the 3D
coordinates of the original structure, with standard deviation of 0.50~\AA. We only use models that
were trained on the wild-type amino acid classification task~\cite{Visani2024-sh}. Models were
pre-trained on structures from 30\% similarity split of ProteinNet's CASP12
set~\cite{alquraishi_proteinnet_2019}, resulting in around $\sim$10k training structures.

We refer the reader to refs.~\cite{Pun2024-qz, Visani2024-sh} for further details on the
architecture, training procedure, training data preparation, and the mathematical introduction to
SO(3)-equivariant models in Fourier space.




\section{TCR-pMHC binding affinity and T-cell activity prediction}
\label{sec:aff_act}

\noindent \textbf{Peptide energy calculation with HERMES.} To predict the affinity of TCR-pMHC
complexes or the activity of TCRs induced by their interactions with a pMHC complex, we introduce
the predicted energy for a peptide $\sigma$, given its surrounding TCR and MHC,

\EQ
E_{\text{peptide}}(\sigma; \text{TCR, MHC}) = \sum_{i=1}^\ell E(\sigma_i ; \text{TCR,
MHC},\sigma_{/i}).
\label{eq:pepEn}\EE

\noindent We assume that each peptide residue contributes linearly to the total energy by the amount
$E(\sigma_i ; \text{TCR, MHC}, \sigma_{/i})$, where $\sigma_i$ is the amino acid at position $i$ of
the peptide. A residue's energy contribution $E$ is evaluated as its associated logit value from
HERMES, taking in as input the atomic composition of the surrounding TCR, MHC and the rest of the
peptide $\sigma_{/i}$ (Fig.~1A).%This method is akin to predicting the effect of a single amino-acid
substitution on a protein's function by using the difference in logits - or equivalently of
log-probabilities - corresponding to the mutant and wild-type amino-acids~\cite{Visani2024-sh}.\\

To compute peptide energies, HERMES operates on the structural composition of the protein complex.
However, data on the impact of substitutions on the TCR-pMHC structure is limited, and computational
models often fail to capture subtle conformational changes in the structure due to amino acid
substitutions in a peptide. To address this issue, we adopt two protocols to estimate energy for a
diverse array of peptides for a given TCR-MHC complex:

\begin{itemize}
\item HERMES-{\em fixed}: In the simplest approach, we select the TCR-pMHC structure with matching
TCR and MHC to our query and with the peptide sequence that most closely matches the peptide of
interest and directly input it into HERMES. The peptide energy is subsequently calculated as
described in eq.~\ref{eq:pepEn} (Fig.~1B). This method does not modify the underlying structure,
effectively assuming that any amino acid substitutions do not significantly alter the protein
conformation.

\item HERMES-{\em relaxed}: Since amino acid substitutions can locally reshape the protein complex,
we introduce a more involved protocol to account for these structural changes. We begin with the
available crystal structure of the TCR-pMHC complex that closely matches our query as template. We
then mutate the original peptide to the desired state and relax the structure using PyRosetta's
cartesian\_ddg fastrelax protocol~\cite{Chaudhury2010-mo}. During relaxation, side chain and
backbone atoms of the peptide are allowed to move, while only the side chain atoms of the MHC and
TCR chains within a 10\AA\, of the peptide are flexible. We then calculate the peptide energy with
HERMES (eq.\ref{eq:pepEn}), using the relaxed structure of the mutated variant as input (Fig.~1B).
Since PyRosetta's relaxation procedure is stochastic in nature, we run 100 realizations of the
relaxation procedure and average the peptide energies across them.
 For ablation purposes, we also use an alternative protocol, whereby we use the HERMES peptide
 energy from the conformation with the minimum (most favorable) Rosetta energy computed on the
 relaxed atoms; we call this protocol {\em relaxed-min\_energy}.
\end{itemize}

\noindent Overall, we use four protocols for scoring peptides that differ in their restrained
protocol (HERMES-{\em fixed} or HERMES-{\em relaxed}), and the amount of noise injected during
training of the underlying HERMES model. Specifically, we train HERMES models with no added noise
($\sigma =0.00~\AA$) and another with moderate Gaussian noise with standard deviation $\sigma=
0.50~\AA$~applied to the atomic coordinates of the protein structure data during training. We refer
to the resulting peptide scoring protocols as HERMES-{\em fixed} 0.00, HERMES-{\em fixed} 0.50,
HERMES-{\em relaxed} 0.00, HERMES-{\em relaxed} 0.50, specifying both the model and the noise
amplitude.




\noindent{\bf Evaluation of EC$_{50}$ from dose response measurements against peptide libraries.}
Ref.~\cite{Luksza2022-gx} presents dose response measurements for fraction of 4-1BB$^+$ CD8$^+$
T-cells across varying peptide concentrations. These measurements are done for peptide libraries of
single point mutants from a given wild-type peptide, across 7 different TCR-MHC systems: three TCRs
recognizing variants of the highly immunogenic HLA-A02:01-restricted human cytomegalovirus (CMV)
epitope NLVPMVATV (NLV); three TCRs recognizing variants of the {weaker} HLA-A02:01-restricted
melanoma gp100 self-antigen IMDQVPFSV; and one TCR recognizing variants of the {weakly immunogenic}
HLA-B*27:05-restricted pancreatic cancer neo-peptide GRLKALCQR~\cite{Luksza2022-gx}.

The half-maximal effective concentration (EC$_{50}$) from the dose response curves of 4-1BB$^+$
CD8$^+$ T-cell fractions can quantify the level of T-cell activity in response to each peptide. We
used our own fits to estimate the EC$_{50}$ associated with T-cell responses. Specifically, we
fitted a Hill function with background to the dose response measurements for each peptide from
ref.~\cite{Luksza2022-gx},
\EQ
f(x) = \frac{(A-A_0) x^n}{x^n + K^n}+ A_0
\label{eq:EC50}\EE
where $f(x)$ is the fraction of 4-1BB$^+$ CD8$^+$ T-cells pulsed by a peptide at concentration $x$,
$A$ is the response amplitude, $A_0$ is the background activation, $n$ is the Hill coefficient, and
$K$ is the half-maximal effective concentration EC$_{50}$. Similar to ref.~\cite{Luksza2022-gx}, we
regularized our fitted Hill functions, whereby the deviation of the amplitude $A$ and the Hill
coefficient $n$ from 1, and the background activation $A_0$ from 0 are penalized. The code to fit
these curves is available in our github repository.\\


\section{De-novo peptide design}
\label{sec:peptide_design}

\noindent {\bf Peptide design algorithm.}
Given a TCR and an MHC, we use HERMES to design de novo peptides that can form a stable TCR-pMHC
complex, and potentially, elicit a T-cell response. Our approach starts with an existing template
TCR-pMHC structure. We refer to the peptide in the template as wild-type; see
Table~\ref{table:design_info} for the templates we use in our analyses. Similar to evaluating the
peptide energy, we use two pipelines to sample novel peptides:

\begin{itemize}
\item HERMES-{\em fixed}: This is our basic pipeline, in which we generate novel peptide sequences
by conditioning on the fixed TCR-pMHC template structure.
Specifically, we sequentially mask the atoms of each amino acid along the peptide and then use
HERMES to compute the probability $P(\sigma_i|{\bf x}_i)$ of different amino acid types $\sigma_i$
at the masked residue $i$, conditioned on the surrounding atomic neighborhood ${\bf x}_i$, within a
10~\AA \, radius of the masked residue's C-$\alpha$. The neighborhood ${\bf x}_i$ can include atoms
from the TCR, the MHC, and the other amino acids in the peptide. Repeating this procedure for every
residue along the peptide allows us to construct a Position Weight Matrix (PWM) that represents the
probabilities of different amino acids at each peptide position, conditioned on the surrounding
structural neighborhood. We generate peptides using multinomial sampling from the PWM; see the
schematic in Fig.~1C.

\item HERMES-{\em relaxed}: This pipeline incorporates simulated annealing and Markov chain Monte
Carlo (MCMC) sampling to account for structural changes resulting from peptide substitutions. We
start with the template TCR-MHC structure and an initial random peptide sequence $\sigma^{(0)}$,
which is {\em in-silico} packed via PyRosetta, followed by the Relax protocol, yielding a relaxed
structure ${\bf x}^{(0)}$. During relaxation, both side chain and backbone atoms of the peptide are
allowed to move, while only the side chain atoms of the MHC and TCR within 10~\AA~of the peptide are
flexible. Following the HERMES-{\em fixed} protocol, we then evaluate the PWM associated with this
relaxed structure and use it to generate a new peptide sequence, $\sigma^{(1)}$. This sequence is
again packed and relaxed, producing a new structure, ${\bf x}^{(1)}$. Using eq.~\ref{eq:pepEn}, we
evaluate the energies of the peptides in their respective (relaxed) structures as $E(\sigma^{(0)} ;
{\bf x}^{(0)})$ and $E(\sigma^{(1)} ; {\bf x}^{(1)})$, and apply the Metropolis-Hastings criterion
to accept the transition from $\sigma^{(0)}$ to $\sigma^{(1)}$ with probability:
\EQ P_{\sigma^{(0)}\to \sigma^{(1)}} = \text{min}\left(1, e^{- \frac{ E(\sigma^{(1)} ; {\bf
x}^{(1)})- E(\sigma^{(0)} ; {\bf x}^{(0)})}{T}}\right)\EE
where $T$ is the temperature. This process is iterated, incrementally reducing $T$ at each step,
until convergence. The peptide obtained at the final iteration represents the protocol’s de novo
design. Repeating this procedure many times produces a diverse set of peptides spanning various
distances from the wild-type sequence; see the schematic in Fig.~1C and
Algorithm~\ref{alg:hermes_relaxed_design} for more details, and the the convergence curves in
Fig.~\ref{fig:joint_design_distributions}C,F,I.\\
\vspace{1cm}

{
\begin{algorithm}[H]
\figcaption{HERMES-{\em relaxed} protocol for generating TCR-MHC-specific peptide sequences.}
{For brevity, we define ${\bf x}_i \equiv (\text{TCR}, \text{MHC}, \sigma_{/i})$ and drop the arguments
of $E_{\text{peptide}}$.}
\label{alg:hermes_relaxed_design}
\textbf{Input:} TCR-pMHC structure ${\bf x}$, maximum number of iterations N\;
\textbf{Output:} peptide sequence $\sigma^{(N)}$\;
$T_0 \gets 10$\;
$T_f \gets 0$\;
$a \gets 0.95$\;
$T^{(0)} \gets T_0$\;
$\sigma^{(0)} \gets \text{random sequence}$\;
${\bf x}^{(0)} \gets \text{RosettaMutate}({\bf x}^{(0)}, \sigma^{(0)})$\;
${\bf x}^{(0)} \gets \text{RosettaFastRelax}({\bf x}^{(0)})$\;
$E^{(0)}_{\text{peptide}} \gets \sum_{i \in \mathcal{R}_{\text{peptide}}} E(\sigma^{(0)}_i ; {\bf
x}^{(0)}_i)$\;
\For{$k = 1, 2, \dots, N$}{
    \For{$i \in \mathcal{R}_{\text{peptide}}$}{
        $\sigma^{(\text{temp})}_i \sim \text{Multinomial}(\text{softmax}([E(\text{aa}; {\bf
        x}^{(k-1)}_i) / T^{(k)}]_{\text{aa} \in \text{amino-acids}}))$\;
    }
    ${\bf x}^{(\text{temp})} \gets \text{RosettaMutate}({\bf x}^{(k-1)}, \sigma^{(\text{temp})})$\;
    ${\bf x}^{(\text{temp})}  \gets \text{RosettaFastRelax}({\bf x}^{(\text{temp})})$\;
    $E^{(\text{temp})} _{\text{peptide}} \gets \sum_{i \in \mathcal{R}_{\text{peptide}}}
    E(\sigma^{(\text{temp})}_i ; {\bf x}^{(\text{temp})}_i)$\;
    $p \sim \text{Uniform}(0, 1)$\;
    \eIf{$p < \exp(-(E^{(\text{temp})}_{\text{peptide}} - E^{(k-1)}_{\text{peptide}}) / T^{(k)})$}{
        $E_{\text{peptide}}^{(k)} \gets E^{(\text{temp})} _{\text{peptide}}$\;
        ${\bf x}^{(k)} \gets {\bf x}^{(\text{temp})}$\;
        $\sigma^{(k)} \gets \sigma^{(\text{temp})}$\;
    }{
        $E_{\text{peptide}}^{(k)} \gets E_{\text{peptide}}^{(k-1)}$\;
        $x^{(k)} \gets {\bf x}^{(k-1)}$\;
        $\sigma^{(k)} \gets \sigma^{(k-1)}$\;
    }
    $T^{(k+1)} \gets (T_0 - T_f)a^k + T_f$\;
}
\end{algorithm}
}
\vspace{1cm}
\end{itemize}


\noindent Similar to scoring, we use four protocols for peptide design that differ in their
restrained protocol (HERMES-{\em fixed} or HERMES-{\em relaxed}), and the amount of noise injected
during training of the underlying HERMES model. Specifically, we train HERMES models with no added
noise ($\sigma =0.00~\AA$) and another with moderate Gaussian noise with standard deviation $\sigma=
0.50~\AA$ applied to the atomic coordinates of the protein structure data during training. We refer
to the resulting peptide design protocols as HERMES-{\em fixed} 0.00, HERMES-{\em fixed} 0.50,
HERMES-{\em relaxed} 0.00, HERMES-{\em relaxed} 0.50, specifying both the model and the noise
amplitude.\\





\noindent{\bf Filtering metrics for peptide design with HERMES and benchmarking of different design
approaches.} For each system, we first selected designs predicted by NetMHCPan to be at least weak
binders to their target MHCs (EL\_rank $\leq 2.0$)~\cite{reynisson_netmhcpan-41_2020}. We then used
TCRdock~\cite{Bradley2023-rl} to filter peptides for reactivity. TCRdock is a structure prediction
algorithm built on top of AlphaFold2~\cite{Jumper2021-hl}, and refined to model TCR-pMHC
interactions. Its PAE score for the TCR-pMHC interface can be used to distinguish reactive from
decoy peptides~\cite{Bradley2023-rl}, and therefore, can provide a metric for filtering our designed
peptides; other work has similarly used the PAE of structure-prediction algorithms to filter
designed protein binders~\cite{Liu2024-vi}. Using TCRdock, we retained those peptides whose PAE
scores at the TCR-pMHC interface fell below a threshold, as illustrated in
Figs.~\ref{fig:design_tcrdock_nyeso},~\ref{fig:design_tcrdock_ebv},
and~\ref{fig:design_tcrdock_mage}. The distribution of PAE scores is highly system
specific~\cite{Liu2024-vi}, so we set each system's threshold near the TCRdock PAE value of its {\em
wild-type} peptides in complex with the respective TCR-MHC, as the wild types are known binders.
Among the filtered peptides, we selected 93 for experimental validation, ensuring that designs from
both HERMES-{\em fixed} and HERMES-{\em relaxed}, as well as from models with 0.0 and 0.5\AA\, noise
levels were represented. We also enforced a minimum difference of 2-3 amino acids (depending on the
system) between the designs and the wild-type peptides.

It should be noted that the AF3 PAE score can also discriminate between reactive and decoy peptides
(Figs.~\ref{fig:design_scatterplots},~\ref{fig:design_roc_curves}), making it also suitable for
filtering. However, a full-access implementation of AF3 was unavailable before our experimental
validations. Therefore, we used AF3 only as a secondary filter for designs in the EBV and MAGE
systems, and for benchmarking against TCRdock's PAE score for the designed peptides
(Figs.~\ref{fig:design_scatterplots}); see details below.\\


\noindent{\bf Design parameters for different TCR-pMHC systems.} Here, we provide the details on the
specifics of peptide selection for each TCR-MHC system. It should be noted that as we progressed
from NY-ESO to EBV and MAGE, we chose to explore the sequence space farther from the wild-type
peptides, leading to more challenging designs.


\begin{itemize}

\item {\em NY-ESO:} The wild-type peptides (NYESO: SLLMWITQC and NYESO V9C: SLLMWITQV) have TCRdock
PAEs of $5.17$ and $5.18$, respectively.
As positive designs, we selected 58 peptides with the lowest TCRdock PAE values (all below 5.35) and
a Hamming distance of at least two amino acids from their respective wild-type sequences. For
negative designs, we uniformly sampled 35 peptides among those designed by HERMES whose TCRdock PAE
values fell between 5.5 and 7.0. We did not enforce all HERMES models to be represented equally
among these designs, nor did we enforce equal representation of the template structures. However, we
note that all models are represented among positive designs
(Fig.~\ref{fig:design_experimental_nyeso}).

\item {\em EBV:} The wild-type peptides (HPVG: HPVGEADYFEY and HPVG E5Q: HPVGQADYFEY) have TCRdock
PAEs of $4.90$ and $4.82$, and AF3 PAEs of $1.15$ and $1.13$, respectively. We selected 93 positive
peptides using a TCRdock's PAE score cutoff of $5.1$ and AF3 PAE score cutoff of $2.0$.
Specifically, we sampled uniformly below these cutoffs with the following design protocols:
HERMES-{\em fixed} at both noise levels using the HPVG structure (3MV7) as template (24 peptide),
HERMES-{\em fixed} at both noise levels, using the HPVG E5Q structure (4PRP) as template (23
designs); HERMES-{\em relaxed} at both noise levels, using the HPVG structure (3MV7) as template (24
peptides), and HERMES-{\em relaxed} at both noise levels using the HPVG E5Q structure (4PRP) as
template (23 designs); see Fig.~\ref{fig:design_experimental_ebv}.


\item {\em MAGE:} The wild-type peptides (MAGE-A3: EVDPIGHLY and Titin: ESDPIVAQY) have TCRdock PAEs
of $5.21$ and $5.03$, and AF3 PAEs of $1.00$ and $0.98$, respectively. We selected 93 positive
peptides using a TCRdock's PAE score cutoff of 5.2; the AF3 PAE scores for all of these designed
peptide are comparable to that of the wild-types. For these designs, we sampled uniformly below the
PAE cutoffs with the following design protocols: HERMES-{\em fixed} at both noise levels using
MAGE-A3 structure (5BRZ) as template (15 peptides); HERMES-{\em fixed} at both noise levels using
the MAGE-A3 structure (5BRZ) as template with E$_1$ fixed (15 peptides); HERMES-{\em relaxed} at
both noise levels using the MAGE-A3 structure (5BRZ) as template (16 peptides); HERMES-{\em relaxed}
at bothsee Fig.~\ref{fig:design_experimental_ebv} noise levels using the MAGE-A3 structure (5BRZ) as
template with E$_1$ fixed (15 peptides); HERMES-{\em fixed} at both noise levels and using the Titin
structure (5BS0) as template (16 peptides); HERMES-{\em relaxed} at both noise levels using the
Titin structure (5BS0) as template (16 peptides); see Fig.~\ref{fig:design_experimental_mage}.
\end{itemize}

\noindent{\bf Characterizing the size of the sequence space for design.} The size of the sequence
space to explore grows exponentially with increasing distance from the wild-type peptide. If no
restriction is imposed (i.e., all amino acids are permissible with equal probabilities), the size of
the sequence space follows,
\EQ
\Omega_0(d) = {\ell \choose d} 19^d
\EE
where $\ell$ is the length of the peptide, and $d$ is the distance from the wild-type sequence.

MHC presentation imposes constraints on amino acid usages in peptides, especially at the anchor
residues. Let $p_i(a)$ denote the probability of amino acid $a$ at position $i$ of the peptide,
e.g., set by the sequence-specific preferences of the MHC allele. The probability of introducing a
mutation at position $i$ depends on the prominence of the wild-type amino acid, and is equal to
$q_i= 1-p_i(\text{wt})$. The typical number of non-wildtype amino acids realized at position $i$ can
be estimated by $ n_i = 2^{S_i}$, where $S_i = -\sum_{a \neq \text{wt}} \hat p_i(a) \log_2 \hat
p_i(a) $ is the entropy associated with the non-wildtype amino acid frequencies, and $\hat p_i(a) =
\frac{p_i(a)}{1-p_i(\text{wt})}$ is the normalized probability of amino acid $a$, obtained by
excluding the wild-type from the set. Therefore, the typical size of the sequence space at Hamming
distance $d$ from the wild-type peptide---subject to constraints imposed by the amino acid
frequencies $p_i(a)$ (e.g., due to MHC restriction)---is given by,
\EQ
\Omega (d) = \sum_{\{i_1,\dots,i_d\}} \left(q_{i_1} n_{i_1} \right)\dots \left(q_{i_d} n_{i_d}
\right) \equiv \sum_{\{i_i,\dots,i_d\}} 2^{ \hat S_{i_1}+ \dots +\hat S_{i_d} } \label{eq:Omega_d}
\EE
where the summation runs over all ${\ell \choose d}$ ways to choose $d$ distinct sites at which
mutations may occur, and $q_i n_i$ is the expected number of substitutions associated with site
$i$---evaluated as the likelihood of mutating site $i$ times the typical number of viable amino acid
replacements at that site---and $\hat S_i =S_i + \log_2 q_i$ is the resulting adjusted entropy at
site $i$.

In principle, evaluating $\Omega (d) $ in eq.~\ref{eq:Omega_d} involves summation over a
combinatorially large number of terms, which can become computationally prohibitive. In practice,
however, many of the sites have the same adjusted entropy, which simplifies the evaluations of
$\Omega (d) $. For example, the adjusted entropies for mutations away from the 9 amino acid NY-ESO
peptide restricted by HLA-A*02:01 yield the following (rounded) values: $[4, 2, 4, 4, 4, 3, 4, 4,
2]$ in bits, indicating six sites with 4 bits of entropy (or $2^4=16$ typical substitutions), one
site with 3 bits (8 typical substitutions), and two sites with 2 bits (4 typical substitutions). To
demonstrate, in this case, the number of sequences at distance $d=5$ from the wild-type follows,

\EQA
\nonumber \Omega_\text{NYESO} (d=5) &=& {6 \choose 5} 16^5 + {6 \choose 4} 16^4\times \left[ {2
\choose 1} 4^1 + 8^1 \right]+
{6 \choose 3} 16^3\times \left[  {2 \choose 1} 4^1\times  8^1 + {2 \choose 2}  4^2    \right]+ 
{6 \choose 2}16^2  \times {2 \choose 2}  4^2\times 8^1  \\
\nonumber&=& 29,065,216\\
\EEA
Similar calculations can be done to characterize the size of the sequence space at varying distances
from the wild-type, as shown in Fig. 4.

\section{Alternative computational methods for scoring peptides and design}
\label{sec:AltModel}
 We compare our approach to alternative methods (listed below), based on which we characterize
 peptide scores akin to HERMES's peptide energy, and when possible, use them to design novel
 peptides. A detailed comparison between the accuracy of different methods in predicting the
 affinity of TCR-pMHC interactions and the peptide-induced activation of T-cells is presented in
 Tables~\ref{table:spearmanr_affinity} and~\ref{table:spearmanr_activity}. Comparison between the
 fidelity of the designs--only according to computational metrics--is shown in
 Figures~\ref{fig:design_tcrdock_nyeso},~\ref{fig:design_tcrdock_ebv},~and~\ref{fig:design_tcrdock_mage}.
 All the designed peptides as well as the code to generate samples are available in our github
 repository.

\begin{itemize}
\item {\em BLOSUM62 substitution matrix}~\cite{Henikoff1992-jq}: The BLOSUM62 substitution matrix
characterizes the favorability of substitutions between all pairs of amino acids, and has been used
widely to characterize the mutational effects in proteins. As our first (and simplest) benchmark, we
used BLOSUM62 to define a peptide energy (score). Specifically, we score a peptide by summing over
the BLOSUM62 scores associated with the amino acid substitutions from the wild-type peptide $
\sigma^{\text{wt}}$ (the one present in the template structure) to the query peptide sequence
$\sigma$ that we are interested in,
\EQ
E^{\rm (B62)}_{\text{peptide}}(\sigma; \sigma^{\text{wt}}) = \sum_{i=1}^\ell
B62[\sigma^{\text{wt}}_i \to \sigma_i]
\label{eq:score_blosum}
\EE
Notably, the BLOSUM62 predicted peptide energy only depends on the substitutions in the peptide
relative to the wild-type amino acid, and does not depend on the structural context of these
substitutions (Fig.~\ref{fig:hsiue_h1_h8}).

For designs with BLOSUM62, we independently sample amino acids along the peptide based on their
favorability with respect to the wild-type amino acid at each position. Specifically, at a given
position $i$ we interpret the substitution score $B62[\sigma^{\text{wt}}_i \to \sigma_i]$ from the
wild-type to any of the other 19 amino acids as the energy difference between the two states, and by
applying a softmax function with temperature $T$, we evaluate the probability of these possible
substitutions. We then use multinomial sampling to draw new amino acids from this probability
distribution at each position,
\EQ
\sigma_i \sim \text{Multinomial}\left(\text{softmax}\left[\frac{B62[\sigma_i^\text{wt} \to
\sigma]}{T}\right]_{{\sigma} \in \text{amino-acids}}\right)
\EE
For our analyses, we use temperatures $T=1.0, 2.0, 3.0$ to modulate the variability of the designed
peptides.\\


\item{\em \L uksza et al. $M$ and $d \cdot M$ substitution matrices}~\cite{Luksza2022-gx}: Here, $M$
is an amino acid substitution matrix, and $d$ is a vector of position-specific scale factors for
9-mer peptides. Briefly, they were constructed to fit the EC$_{50}$ measurements on seven saturation
mutagenesis experiments of three different wild-type peptides, against a total of seven TCR-MHC
pairs--the very same T-cell activity data that we used to test HERMES in Fig.~3. Specifically, $d$
and $M$ were constructed such that:
\EQ
d_i \cdot M[\sigma^{\text{wt}}_i \to \sigma^{\text{mt}}_i] \approx \log
\frac{\text{EC}_{50}^{\text{mt}}}{\text{EC}_{50}^{\text{wt}}} + R_\text{B62}[\sigma^{\text{wt}}_i
\to \sigma^{\text{mt}}_i]
\label{eq:luksza_matrix}
\EE
where the mutant and wild-type peptides only differ at position $i$, and $R_\text{B62}$ is a
regularization term reflecting the BLOSUM62 substitution matrix; we refer the reader to the
Supplementary Methods of ref.~\cite{Luksza2022-gx} for details. Since all peptides considered when
constructing $d$ and $M$ are 9-mers, the position scale factors can only be used on 9-mer peptides.
We use $d$ and $M$ for scoring in a manner analogous to BLOSUM62:
\EQ
E^{(d \cdot M)}_{\text{peptide}}(\sigma; \sigma^{\text{wt}}) = \sum_{i=1}^\ell d_i \cdot
M[\sigma^{\text{wt}}_i \to \sigma_i]
\label{eq:score_luksza_d_m}
\EE
and
\EQ
E^{(M)}_{\text{peptide}}(\sigma; \sigma^{\text{wt}}) = \sum_{i=1}^\ell M[\sigma^{\text{wt}}_i \to
\sigma_i]
\label{eq:score_luksza_m}
\EE
Similar to BLOSUM62, the predicted peptide energies $E^{(d \cdot M)}_{\text{peptide}}$ and
$E^{(M)}_{\text{peptide}}$ only depend on the type substitutions, and not on their structural
contexts.\\



\item {\em MHC position weight matrices}~\cite{Tadros2023-ij}: MHC alleles exhibit distinct
sequence-specific preferences for peptide presentation, most notably in the peptide's amino acid
composition at anchor positions. Computational models have characterized these preferences with
position weight matrices (PWMs) that reflect the
the probability of different amino acids at each position for peptides presented by a given MHC
allele~\cite{Lundegaard2008-ao, Lundegaard2008-wz, Andreatta2016-uh,Tadros2023-ij}. As a benchmark
for our designs, we use the PWMs from the MHC Motif Atlas~\cite{Tadros2023-ij} to independently
sample amino acids at each site along the peptide. The resulting peptide sequences should reflect
biases imposed by MHC presentation but not interactions with TCRs.\\

\item {\em ProteinMPNN}~\cite{Dauparas2022-fu}: ProteinMPNN is a commonly used tool for
inverse-folding (i.e., designing protein sequences consistent with the backbone of a structure). It
is primarily used to sample amino acid sequences conditioned on a protein's backbone structure, and
optionally with partial sequence information (i.e., amino acids specified at some of the residues).
As ProteinMPNN also outputs probability distributions of amino acids at different sites of the
protein, it can also be used to score peptides (i.e., evaluate peptide energy) in a manner similar
to HERMES, as

\EQA
 E^{\rm (MPNN)}_{\text{pep 1-by-1}}(\sigma; TCR, MHC) &=& \sum_{i=1}^\ell \log P^{\rm
 (MPNN)}(\sigma_i ; {\rm backbone, \sigma_{\text{TCR}}, \sigma_\text{MHC}}, \sigma_{/i})
\label{eq:peplogp_proteinmpnn}
\EEA
where $P^{\rm (MPNN)}(\sigma_i ; {\rm backbone, TCR, MHC}, \sigma_{/i})$ is the ProteinMPNN estimate
for the probability of amino acid $\sigma_i$ at position $i$ of the peptide, given the backbone
structure of the whole protein complex, and the amino acid sequence of the MHC $\sigma_\text{MHC}$,
TCR $\sigma_\text{TCR}$, and the rest of the peptide $\sigma_{/i}$. We term this scoring scheme as
peptide 1-by-1 since we mask the peptide's residues one at a time.

We explored two additional ProteinMPNN scoring schemes that differ in the residues they mask and
score. As our second scoring scheme, we simultaneously mask the entire peptide's sequence $\sigma$,
instead of one amino-acid at a time, and evaluate the score for the whole peptide $E^{\rm
(MPNN)}_{\text{pep}}(\sigma; TCR, MHC)$:
\EQA
 E^{\rm (MPNN)}_{\text{pep}}(\sigma; TCR, MHC) &=& \sum_{i=1}^\ell \log P^{\rm (MPNN)}(\sigma_i ;
 {\rm backbone, \sigma_{\text{TCR}}, \sigma_\text{MHC}})
\label{eq:peplogp_proteinmpnn_pep}
\EEA

For the third scoring scheme, we consider masking and scoring the TCR residues that make contact
with the peptide--defined as those whose C-$\alpha$ lie within 12~\AA~ of any peptide residue's
C-$\alpha$. We call this set {\bf TCR-int} and its corresponding amino acid composition as $\sigma
_\text{TCR-int}$. In this scheme, both the peptide and the interacting TCR residues are masked
simultaneously and scored together:
\EQA
\nonumber E^{\rm (MPNN)}_{\text{pep and tcr}}(\sigma; TCR, MHC) &=& \sum_{i=1}^\ell \log P^{\rm
(MPNN)}(\sigma_i ; {\rm backbone, \sigma_{\text{TCR / TCR-int}}, \sigma_\text{MHC}}) \\ &+&
\sum_{j=1}^n \log P^{\rm (MPNN)}(\sigma_{\text{TCR-int}_j} ; {\rm backbone, \sigma_{\text{TCR /
TCR-int}}, \sigma_\text{MHC}})
\label{eq:peplogp_proteinmpnn_tcr_and_pep}
\EEA

Similar to HERMES, we considered ProteinMPNN models trained with two noise levels: 0.02~\AA~(very
little noise) and 0.20 \AA.

For designs with ProteinMPNN, we use the algorithm's native functionalities, keeping the TCR and MHC
structures fixed (from the design template) and prompting the model to design only the peptide. We
use two ProteinMPNN models (trained with 0.02~\AA~and 0.20~\AA~noise), and use a temperature of 0.7
for all systems. \\


\item {\em ESM-IF1}~\cite{Hsu2022-km}: ESM-IF1, like ProteinMPNN, is an inverse-folding model, but
it lacks native support for scoring or designing proteins when only partial sequence information is
available.

To access the ESM-IF1 model, we used the code in the esm repository
at~\url{https://github.com/facebookresearch/esm/tree/main/examples/inverse_folding}.

For scoring a peptide sequence given a TCR-MHC context, we used the script
\texttt{score\_log\_likelihoods.py} with the \texttt{--multichain-backbone} flag, with the same
TCR-pMHC structure inputs as for HERMES and ProteinMPNN. This script computes a score for the
peptide sequence conditioned on the backbone coordinates of the entire TCR-pMHC complex, but {\em
not} on the sequence of the TCR and the MHC:
\EQA
\nonumber E^{\rm (ESM-IF1)}_{\text{peptide}}(\sigma; TCR, MHC) &=& \log P^{\rm
(ESM-IF1)}_{\text{peptide}}(\sigma ; {\rm backbone})
\label{eq:peplogp_esmif}
\EEA

To design peptide sequences given a TCR-MHC context, we used the script
\texttt{sample\_sequences.py} with the \texttt{--multichain-backbone} flag, and a temperature of 0.7
for all systems, with the same TCR-pMHC structure inputs as for HERMES and ProteinMPNN. Similar to
scoring, this script samples peptide sequences conditioned on the backbone coordinates of the entire
TCR-pMHC complex, but {\em not} on the sequences of the TCR and MHC.\\


\item {\em TCRdock}~\cite{Bradley2023-rl}: TCRdock is a structure prediction algorithm built on top
of AlphaFold2~\cite{Jumper2021-hl}, and refined to predict the structure of TCR-pMHC complexes. The
algorithm model the structure of TCR-pMHC complexes by leveraging structural templates provided by
its curated database, as well as the knowledge of TCR-pMHC canonical binding poses. Crucially,
TCRdock's Predicted Alignment Error (PAE) for TCR-pMHC interfaces was found to well discriminate
between true peptide binders and random decoys~\cite{Bradley2023-rl}. Thus, to benchmark against
TCRdock, we interpret this PAE score as peptide energy, and use it to predict TCR-pMHC binding
affinities and peptide-induced T-cell activities.

TCRdock takes as input sequences of a TCR, MHC, and peptide, and identifies TCR-pMHC structures in
its database with similar amino acid sequences to the query to serve as templates to model the
structure of the input TCR-pMHC complex. To examine how template-query sequence similarity
influences TCRdock's predictive accuracy, we ran the algorithm under two conditions: the default
mode, which uses any closest-matching templates available, and the {\em benchmark} mode, which
restricts the maximum sequence similarity of templates with the input TCR-pMHC.

We use TCRdock to characterize peptide energies, and to filter the sequences designed by HERMES.
However, TCRdock does not provide a natural way to sample peptides for design, and because of its
high computational cost, we did not include TCRdock PAE calculations in our design pipeline (e.g.,
for MCMC sampling).\\

\item {\em TULIP}~\cite{Meynard-Piganeau2024-po}: TULIP is a transformer-based model trained to
generate either TCR or peptide sequences conditioned on the other binding partners in a TCR-pMHC
complex. The conditional likelihood function learned by TULIP can be used to score peptides within a
TCR-pMHC context. Specifically, this scoring scheme was used in ref.~\cite{Meynard-Piganeau2024-po}
to predict the T-cell activity measurements from ref.~\cite{Luksza2022-gx}, and we benchmarked our
predictions against TULIP's for this data. In Table~\ref{table:spearmanr_activity}, we present the
Spearman correlations between TULIP predictions and the fitted EC$_{50}$ values as reported in
ref.~\cite{Luksza2022-gx}, which are the same correlations reported in the TULIP
paper~\cite{Meynard-Piganeau2024-po}. Additionally, we provide correlations with the fitted
EC$_{50}$ values obtained using the procedure outlined in Section~\ref{sec:aff_act} and
eq.~\ref{eq:EC50}, consistent with correlations reported in Fig.~3.\\


\item {\em NetTCR2.2}~\cite{Jensen2024-xs}: NetTCR2.2 consists of a Convolutional Neural Network
classifier to predict TCR-peptide binding from the the amino acid sequences of the peptide and of
all six CDRs from the $\alpha$ and the $\beta$ chains of a TCR. Because the training set contains
many examples of different TCRs binding to the same peptide, but far fewer cases of multiple
peptides tested against a given TCR, the model performs best when ranking candidate TCRs for a fixed
peptide, as reported in ref.~\cite{Jensen2024-xs}. By contrast, we aim to address the complementary
task for ranking peptides for a fixed TCR.

We query NetTCR2.2 through the authors' web server
(\url{https://services.healthtech.dtu.dk/services/NetTCR-2.2/}). CDR boundaries are assigned by
aligning each TCR sequence to the IMGT reference sequences, using code from the TCRdock
repository~\cite{Bradley2023-rl}. Of the TCR-pMHC pairs examined in this work, only one system
(TCR1, Table~\ref{table:affinity_and_activity_info}) appears in the NetTCR2.2 training set, and even
then with just a single annotated positive peptide. Therefore, this small overlap is unlikely to
bias our predictions.\\

\item {\em TAPIR}~\cite{Fast2023-tk}: Similar to NetTCR-2.2, TAPIR consists of a Convolutional
Neural Network trained to predict TCR-peptide-MHC binding. As input, TAPIR considers the CDR3's
amino acid sequences, the V and J genes, the peptide's amino acid sequence, and the MHC allele. We
use TAPIR for scoring via the web interface hosted by the authors~\url{https://vcreate.io/tapir/}.
We extract V/J genes and CDR3 from a TCR sequence by aligning it to the IMGT reference sequences,
using code from the TCRdock repository~\cite{Bradley2023-rl}. We find that five out of nine of the
TCR-pMHC systems in this study are present in the training data of TAPIR, but each only with a
single annotated positive peptide (1G4 TCR, TCR1, TCR4, TCR5, TCR6,
Table~\ref{table:affinity_and_activity_info}), and therefore, this small overlap is unlikely to bias
our predictions.

\end{itemize}


\noindent \textbf{Computational modeling of TCR-pMHC templates with AlphaFold3.} For six out of
seven systems with T-cell activity measurements from ref.~\cite{Luksza2022-gx} (TCRs~2-7), no
experimentally-determined structure is available. For these, we generated structures using the
AlphaFold3 (AF3) model~\cite{Abramson2024-jp}. Furthermore, to probe how sensitive HERMES is to AF3
modeled inputs, we produced AF3 predictions for every TCR-pMHC complex we studies in this work--both
for scoring and design design targets--using the wild-type sequences of all the chains as inputs;
see Dataset~\ref{dataset:af3_sequences} for these sequences. Each target was folded using both the
AF3 server's option with template search (with default cutoff date of September 9th 2021), and
without template, the latter approximating the scenario in which no homologous structures exist. We
note, however, that several of TCR-pMHC complexes we analyzed here were likely present in the AF3
training data, so the ``template-free" predictions can only serve as an approximate test of true
de-novo performance.

We fold TCR-pMHC systems for which we have a structure using restricted amino-acid sequences of the
TCR and MHC chains that were extracted with functionalities in the TCRdock codebase, and are
equivalent to the sequences that TCRdock would use to predict the structure~\cite{Bradley2023-rl}.
Specifically, we first extract the amino-acid sequences from the experimental structure. Then, we
extract V/J genes, CDR3 and MHC allele from them using `parse\_tcr\_pmhc\_pdbfile.py`, and input
them to `setup\_for\_alphafold.py' to get the restricted amino-acid sequences. For the bispecific
Fab antibody-mediated T-cell therapy from ref.~\cite{Hsiue2021-ag}, we used AF3 to fold and dock the
Fab and the pMHC by extracting the amino-acid sequences from the experimental structure. We found
that AF3 displayed high uncertainty (as measured by plDDt) only at the Fab-pMHC interface, and not
in the folds of individual chains. The sequences we used as input to AF3 are listed in
Dataset~\ref{dataset:af3_sequences}. \\
\\


\noindent \textbf{Benchmarking HERMES predictions for T-cell affinity and activity.} We compared the
performance of HERMES models at predicting the binding affinity and activity of T-cells against
substitution matrices (BLOSUM62~\cite{Henikoff1992-jq}, and \L uksza {\em et al}'s TCR-specific $d
\cdot M$ and $M$ matrices~\cite{Luksza2022-gx}), the TCR-pMHC structure prediction algorithm
TCRdock~\cite{Bradley2023-rl}, the sequence-based machine learning algorithms for TCRs
(TAPIR~\cite{Fast2023-tk}, NetTCR~\cite{Jensen2024-xs}, and TULIP~\cite{Meynard-Piganeau2024-po}),
and the structure-based inverse folding models (ESM-IF1~\cite{Hsu2022-km} and
ProteinMPNN~\cite{Dauparas2022-fu}). Since all of these models are unsupervised in nature (except
for \L uksza {\em et al} activity predictions for TCRs 1-7), we do not expect predictions to be in
the same units as the affinity and activity values, and thus, we use Spearman correlation between
data and predictions as a metric for model performance.

For predicting TCR-pMHC binding affinities (Table~\ref{table:spearmanr_affinity}), 
BLOSUM62 leads for the 1G4 TCR system, while template-guided TCRdock tops on the A6 TCR systems.
HERMES-{\em relaxed} with no noise (0.00) achieves slightly lower Spearman correlation but is
statistically indistinguishable from these best performers (p-value $>0.05$ associated with the
difference in the Fisher's z-transformed Spearman correlations). For 1G4, ESM-IF1 and the \L uksza
{\em et al}'s TCR-specific $d \cdot M$ and $M$ substitution matrices are similarly competitive.
HERMES-{\em relaxed} 0.00 consistently outperforms the sequence-based models (TAPIR, NetTCR-2.2),
ProteinMPNN, and TCRdock-{\em no template} (the TCRdock model without using structural templates
similar to the query data). Lastly, choosing the minimum-Rosetta-energy conformation for HERMES-{\em
relaxed} yields correlations with the data that are comparable to the baseline HERMES-{\em relaxed}
model that averages over the energies of 100 different realizations of relaxed peptides (Fisher's
z-test p-value $>0.05$). Table~\ref{table:spearmanr_affinity} shows these benchmarking results in
detail.


For predicting T-cell activity (Table~\ref{table:spearmanr_activity}), HERMES-{\em fixed} without
noise (0.00) achieves the best performance for TCR1, TCR2, and TCR3, and proteinMPNN-0.02 shows
competitive performances for TCR1 and TCR3. ProteinMPNN-0.02 leads for TCR6 and H2-scDb, while
HERMES-{\em fixed} 0.00 achieves slightly lower but statistically indistinguishable Spearman
correlations (Fisher's z-test p-value $>0.05$). All models struggle to provide reliable predictions
for TCR7, with the exception of ProteinMPNN 0.02, which attains a moderate Spearman correlation
$\rho=0.31$ when scoring one peptide amino acid at a time (eq.~\ref{eq:peplogp_proteinmpnn}), and
$\rho=0.29$, when scoring the peptide and the TCR together
(eq.~\ref{eq:peplogp_proteinmpnn_tcr_and_pep}). For TCR4, which has a poor structural model (AF3 PAE
= 2.07), all structure-based models fail to produce reliable predictions, whereas BLOSUM62 achieves
a moderate Spearman correlation $\rho=0.33$. The \L uksza {\em et al}'s $d \cdot M$ substitution
matrix overall performs well on the TCR~1-7 benchmark; however, because it was obtained by fitting
on these same data, its predictions are not strictly zero-shot. The structure-based models HERMES
and ProteinMPNN consistently outperform the sequence-based models (TAPIR, NetTCR-2.2, and TULIP), as
well as the structure-based ESM-IF1. Lastly, choosing the minimum-Rosetta-energy conformation for
HERMES-{\em relaxed} model yields a similar performance to the baseline HERMES-{\em relaxed} model
that averages over the energies of 100 different realizations of relaxed peptides.
Table~\ref{table:spearmanr_activity} shows these benchmarking results in detail.

Our results do not indicate a clear preference between HERMES-{\em fixed} and HERMES-{\em relaxed}:
HERMES-{\em fixed} predicted T-cell activity better, while HERMES-{\em relaxed} predicted binding
affinity between TCR and pMHC better. However, it is important to note that in all systems tested
here, the scored peptides are highly similar to the wild-type peptides, with a maximum Hamming
distance of just one amino acid from a wild-type peptide found in one of the available structural
templates. As a result, relaxation may not be essential for achieving high performance.
Interestingly, noised HERMES models underperform compared to the un-noised counterparts. This is in
contrast to prior observations, where noised models achieve higher performances in predicting the
stability effect of mutations in proteins~\cite{Visani2024-sh}.
 Moreover, we observe that ProteinMPNN predictions are robust to the choice of scoring scheme, i.e.,
 scoring peptide amino acids one at a time $E^{\rm (MPNN)}_{\text{pep 1-by-1}}$
 (eq.~\ref{eq:peplogp_proteinmpnn}), the whole peptide simultaneously $E^{\rm (MPNN)}_{\text{pep}}$
 (eq.~\ref{eq:peplogp_proteinmpnn_pep}), or peptide and the TCR interface together $E^{\rm
 (MPNN)}_{\text{pep and tcr}}$ (eq.~\ref{eq:peplogp_proteinmpnn_tcr_and_pep}); see
 Fig.~\ref{fig:proteinmpnn_scoring_schemes_comparison}.

Lastly, we benchmarked all four HERMES models and two ProteinMPNN models (using peptide-scoring
schemes eq.~\ref{eq:peplogp_proteinmpnn_pep} with different noise levels) on AF3-predicted
structures, for both activity and affinity predictions (Fig.~\ref{fig:plots_af3_struc_ablation}). We
generated the AF3 structures both with and without template guidance, and the resulting predictions
made with these AF3 structures were compared to performances using the corresponding experimental
structures. AF3 produces accurate models for all complexes except the H2-scDb Fab bound to pMHC, for
which the docking orientation and conformation were incorrect
(Table~\ref{table:af3_folding_metrics}). Expectantly, we saw a significant drop in model
performances for both HERMES and ProteinMPNN when using the inaccurate AF3 predictions for H2-scDb
(Fig.~\ref{fig:plots_af3_struc_ablation}). For both HERMES and ProteinMPNN, replacing experimental
structures with AF3 predictions altered predictions for the 1G4 and A6 TCRs binding to pMHC
variants, but the changes were not statistically significant (p-value$>0.05$ associated with the
difference in the Fisher's z-transformed Spearman correlations); see
Fig.~\ref{fig:plots_af3_struc_ablation}. Finally, for TCR1, we observe no differences in performance
at predicting T-cell activity between using the AF3 structures vs. the experimental structures.
Comparisons for TCRs~2-7 show insignificant differences between predictions using AF3 models with
and without templates.

In Table~\ref{table:af3_folding_metrics} we report, where available, the C$_\alpha$ RMSD between
AF3-predicted and experimentally determined structures, together with AF3 confidence metrics
computed with and without template guidance. The confidence metrics include the TCR–pMHC interface
predicted alignment error (PAE), the interchain predicted TM-score (ipTM), and the predicted
TM-score (pTM; not interface-specific). We found ipTM and pTM to be closely related
(Table~\ref{table:af3_folding_metrics}). Because structure-based models such as HERMES and
ProteinMPNN rely on accurate input structures, Fig.~\ref{fig:plots_af3_struc_ablation}B plots their
best-performing variants against AF3 confidence. The AF3 model for TCR4—on which both HERMES and
ProteinMPNN perform poorly—exhibits the second-lowest ipTM confidence (after H2-scDb). However, the
limited range of available AF3 confidence values prevents definitive conclusions about how these
metrics impact
the accuracy of predictions for peptide mutation-effects on T-cell responses.
\\

%\GMV{Finally, we tested whether AF3 confidence metrics could be used as a quantitative estimate of
performance for the structure-based HERMES and ProteinMPNN models. Specifically, we hypothesized
that predictions made on low-confidence structures should be expected to be of poor quality, since
structure-based models assume the use of accurate crystal-structure-like structures. To test this,
we plotted the models' performance on AF3-folded structures as a function of AF3's confidence
metrics (Fig.~\ref{fig:plots_af3_struc_ablation}). We considered PAE and interchain predicted
Template Modeling (ipTM) scores, two functionally different scores describing the models' confidence
in its prediction of the structure at the interface between chains. While PAE was computed
specifically on the interface between TCR/Ab and pMHC, we considered the overall ipTM score provided
by AF3 to increase the amount of complementary information between the two metrics. We also report
the predicted Template Modeling (pTM) scores - which are not specific to interfaces - in
Table~\ref{table:af3_folding_metrics}, and find them to be very closely correlated with ipTM scores.
Unfortunately, lack of datapoints with very low confidence metrics prevents the drawing of any
significant conclusions. In fact, all systems considered except for the H2-scDb Fab and TCR4 have an
ipTM score above 0.80 - with TCR4 being right at the boundary - a value that is generally considered
to be the threshold for high-quality predictions~\cite{Abramson2024-jp}.



\noindent \textbf{Benchmarking HERMES designs against alternative models.} Although we did not
experimentally test the designs by the different models, we used both peptide presentation by MHC
predicted by NetMHCPan and TCRdock’s PAE as benchmarking metrics across different design methods. We
note that these are only computational metrics and should not be taken as ground truth; this is
particularly true for results using TCRdock PAE as a proxy for binding, which are much more likely
farther from the ground truth than the easier-to-predict MHC presentation.

Figs.~\ref{fig:design_tcrdock_nyeso}, \ref{fig:design_tcrdock_ebv}, and
\ref{fig:design_tcrdock_mage} show the fraction of peptides meeting each criterion for designs
produced by various HERMES models, ProteinMPNN, BLOSUM62, and the system's MHC peptide preferences.
Consistently, across the three systems, and for all models, the fraction of designs passing the PAE
threshold decreases with increasing Hamming distance from the wild-types' sequences
(Figs.~\ref{fig:design_tcrdock_nyeso},~\ref{fig:design_tcrdock_ebv},
and~\ref{fig:design_tcrdock_mage}; panels E, and G). Also consistently, but more surprisingly,
models trained with Gaussian noise applied to the structure---for both HERMES and
ProteinMPNN---perform worst than their noise-less counterparts
(Figs.~\ref{fig:design_tcrdock_nyeso},~\ref{fig:design_tcrdock_ebv},
and~\ref{fig:design_tcrdock_mage}; panels E, and G).

HERMES models are consistently the best at designing peptides that are scored well for presentation
by the respective MHC alleles (panel C of
Figs.~\ref{fig:design_tcrdock_nyeso},~\ref{fig:design_tcrdock_ebv},
and~\ref{fig:design_tcrdock_mage}). No single model stands out as being consistently best across the
three systems at designing peptides above the TCRdock PAE threshold: HERMES outperforms ProteinMPNN
in the NY-ESO system (Fig.~\ref{fig:design_tcrdock_nyeso}), performs similarly in the EBV system
(Fig.~\ref{fig:design_tcrdock_ebv}), and performs worse in the MAGE system
(Fig.~\ref{fig:design_tcrdock_mage}). BLOSUM62 makes worse designs compared to HERMES and
ProteinMPNN based on TCRdock PAE scores. Surprisingly, ESM-IF1 designs have the lowest success in
MHC presentation, but show better TCRdock PAE scores compared to HERMES-{\em relaxed} for the MAGE
system (Figs.~\ref{fig:design_tcrdock_nyeso},~\ref{fig:design_tcrdock_ebv},
and~\ref{fig:design_tcrdock_mage}). Designs sampled from the MHC motif, while scoring well on MHC
presentation, consistently score poorly for TCRdock PAE. We expected HERMES-{\em relaxed} to design
better binders than HERMES-{\em fixed} at greater sequence distances from the wild-types. However,
this holds only for the NY-ESO system (Fig.~\ref{fig:design_tcrdock_nyeso}) and not for EBV or MAGE
(Figs.~\ref{fig:design_tcrdock_ebv},~\ref{fig:design_tcrdock_mage}). Overall, a broader benchmarking
study, ideally with experimental validation, is needed to rigorously evaluate the effectiveness of
different peptide design methods.

To assess how using AF3 structures as templates influence design, we ran the HERMES-{\em fixed}
design protocol using experimental and AF3 structures as templates and compared the resulting
peptide likelihood matrices (as position weight matrices or PWMs) for the three systems
(Fig.~\ref{fig:joint_logoplots_for_design}). Across NY-ESO, EBV, and MAGE, the resulting PWMs were
largely consistent across the two template choices: most peptide positions showed the same
preferences, and when variability did appear it was typically confined to chemically similar
residues. There are two noteworthy exceptions: in NY-ESO, the top choice in position 7 of the
peptide was flipped from Threonine and Aspartic Acid when AF3 template was used
(Fig.~\ref{fig:joint_logoplots_for_design}A). In MAGE, using the AF3 model introduced a pronounced
preference for glutamic acid at position 1 (Fig.~\ref{fig:joint_logoplots_for_design}C), a residue
that we have already shown to be important for eliciting T-cell activity (Fig.~4). Aside from these
cases, AF3-derived structures preserve the peptide specificity landscape learned by HERMES.


\section{Experimental protocol to measure peptide activities}
\noindent {\bf Artificial antigen-presenting cell lines.} We obtained the full-length coding
sequences of HLA-A*02:01, HLA-A*01:01, and HLA-B*35:01 from the IMGT/HLA
database~\cite{Robinson2003-ru}. Gene fragments were synthesized by Genscript and were cloned into
the lentiviral pLVX-EF1$\alpha$-IRES-Puro (Clontech) vector. We used Lipofectamine 3000 Transfection
Reagent (Invitrogen) to transduce the 293T packaging cell line (ATCC CRL-3216) with the psPAX2
packaging plasmid (Addgene plasmid \#12260), the pMD2.G envelope plasmid (Addgene plasmid \#12259),
and the HLA-encoding lentiviral vector. The lentivirus-containing supernatant was collected after 24
and 48 hours post-transfection and was filtered through a 0.45um SFCA filter (Thermo Scientific).
The lentivirus was concentrated using Lenti-X Concentrator (Takara Bio) according to the
manufacturer’s protocol, resuspended in DPBS (Gibco), and stored at -80C for future use. To generate
monoallelic HLA cell lines for HLA-A*02:01 and HLA-B*35:01, HLA-deficient K562 cells (ATCC CCL-243)
were transduced with lentivirus. Since K562 cells have a high expression of both MAGE-A3 and Titin
antigens, we transduced murine EL4 cell lines (ATCC TIB-39) with HLA-A*01:01 lentivirus to produce
HLA-A*01:01 presenters lacking the natural expression of the tested antigens. At 48 hours
post-transduction, K562 cells were transferred to Iscove’s Modified Dulbecco’s Medium containing
10\% FBS, 2mM L-glutamine, 100 U/mL penicillin/streptomycin, and 2 $\mu$g/mL puromycin (Gibco). EL4
cells were transferred to RPMI media containing 10\% FBS, 2mM L-glutamine, 100 U/mL
penicillin/streptomycin, and 2 $\mu$g/mL puromycin (Gibco). The antibiotic selection continued for
one week before transferring cells to the antibiotic-free media.\\

\noindent {\bf TCR-expressing Jurkat 76.7 cell lines.} Nucleotide sequences of TCRa and TCRb chains
were generated using stitchr~\cite{Heather2022-nr}. Both chains were modified to incorporate murine
constant segments (TRAC*01 and TRBC2*01) to enhance TCR surface expression. For NYESO and
EBV-specific TCRs, the TCRa and TCRb chains were connected through a P2A ribosomal skip motif,
synthesized as a single bicistronic gene fragment by Genscript, and cloned into the lentiviral
pLVX-EF1a-P2A-mCherry-IRES-Puro (Clontech). For MAGE-specific TCR, gene fragments for the TCRa and
TCRb chains were separately synthesized by Genscript and cloned into either the lentiviral
pLVX-EF1a-P2A-mCherry-IRES-Puro (Clontech) or pLVX-EF1$\alpha$-IRES-G418 (modified in-house from
pLVX-EF1$\alpha$-IRES-Puro). The 293T packaging cell line (ATCC CRL-3216) was transfected with
TCR-encoding lentiviral vectors, psPAX2 packaging plasmid (Addgene plasmid \#12260), and pMD2.G
envelope plasmid (Addgene plasmid \#12259) using Lipofectamine 3000 Transfection Reagent
(Invitrogen). The lentivirus-containing supernatant was collected after 24 and 48-hours after
transfection, filtered through a 0.45$\mu$m SFCA filter (Thermo Scientific), concentrated using
Lenti-X Concentrator (Takara Bio) according to the manufacturer’s protocol, resuspended in DPBS
(Gibco), and stored at -80C for future use. To generate Jurkat cells expressing the paired TCR,
TCR-null CD8-positive Jurkat 76.7 cells with NFAT-eGFP reporter were transduced with lentivirus. At
48 hours post-transduction, cells were transferred to RPMI media containing 10\% FBS, 2mM
L-glutamine, 100 U/mL penicillin/streptomycin, and 1 $\mu$g/mL puromycin (Gibco). For MAGE-specific
cell line generation Jurkat cells were co-transduced with lentiviruses encoding TCRa and TCRb chains
and cultured with both 1 $\mu$g/mL puromycin and 500 $\mu$g/mL Genectin (Gibco). The antibiotic
selection continued for one week. The abTCR expression on the surface was confirmed by mCherry
expression and staining with a 1:100 dilution of anti-mouse TCRb APC-Fire750 (clone H57-597, 109246,
Biolegend).\\

\noindent {\bf Specificity validation of TCR-expressing Jurkat.} Original cognate peptides and their
variants were synthesized as a crude peptide library with by Genscript, and then diluted to a 4mM
stock solution in DMSO (Sigma). To examine the specificity of the TCR-expressing Jurkat cell lines,
we co-cultured 100,000 Jurkat cells with equal amounts of artificial APCs expressing the
corresponding HLA allele in a 96-well round-bottom tissue culture plate. We tested each
specificity—MAGE, EBV, NYESO—under 96 different conditions. These included one or two positive
control original peptides and an unstimulated “no peptide” control. The cells were cultured in RPMI
media supplemented with 10\% FBS, 2mM L-glutamine, 100 U/mL penicillin/streptomycin, and 1 $\mu$g/ml
each of anti-human CD28 (BD Biosciences, 555725) and CD49d (BD Biosciences, 555501). The co-cultures
were pulsed with 1$\mu$M peptide in triplicate. We added DMSO to unstimulated control wells (no
peptide) at concentrations equivalent to those in the peptide-stimulated wells. After incubating the
cells for 24 hours at 37$^\circ$C and 5\% CO2, we washed them once in DPBS (Gibco) and resuspended
them in 150$\mu$l DPBS for flow cytometry analysis using a custom-configured BD Fortessa with
FACSDiva software (Beckton Dickenson). We determined GFP percentages using FlowJo version 10.8.1
software (BD Biosciences). Alternatively, we measured the GFP percentage in wells of a 96-well
round-bottom plate using Cellcyte X (Echo), using 4x objective in spheroid mode, acquisition in
green channel (800 ms exposure 3dB gain). Image analysis was done with CellCyte Studio software with
total spheroid area recipe and standard defaults (50 a.u. Contrast Sensitivity, 2 a.u Smoothing, 100
$\mu$m$^2$ filled hole size 100 $\mu$m Min. object size), relative GFP+ spheroid area was used as
output. For NY-ESO, we compared the T-cell activity level measured using flow cytometry with
fluorescence microscopy, and show consistency between the two (Fig.~\ref{fig:flow_vs_microscopy}).
Given the consistency between the two experimental approaches, we relied on fluorescence microscopy
only to measure GFP levels induced by different peptides for MAGE and EBV. \\

\noindent {\bf Selecting the GFP expression threshold for T-cell activity.} Experimental
measurements are done across three replicates, and the reported GFP levels in
Figs.~4,~\ref{fig:design_scatterplots} for each condition are averaged over the three replicates. We
consider a peptide to have activated T-cells if its induced replicate-averaged GFP level met two
criteria: 1) it exceeded $0.5\%$, and 2) it surpassed the mean plus three standard deviations of the
GFP level negative controls (no-peptide).
According to these criteria, the GFP activation cutoffs were 0.5\% for the NY-ESO system, 0.64\% for
the EBV system, and 0.76\% for the MAGE system.



\section{Experimental activities of the designed peptides}
\label{sec:exp_analysis}

\noindent {\bf Comparing design success rates across the three TCR-pMHC systems.} Among the designed
peptides, 50\% in the NY-ESO system (Fig. 4A), 16\% in the EBV system (Fig.~4B), and 26\% in the
MAGE system induce significant T-cell activities. However, it should be noted that the designed
NY-ESO peptides were closer in sequence to their wild-type templates (2-7 amino acid differences)
than the peptides designed for the EBV (3-9 differences) or MAGE (3-8 differences) systems. We find
that the activity of designs decays with increasing distance from the respective wild-type, with
almost no designs beyond 5 amino-acid distance from the wild-types activating their respective
T-cells
(Fig.~4,~\ref{fig:design_experimental_nyeso},~\ref{fig:design_experimental_ebv},~\ref{fig:design_experimental_mage}).
This underscores the difficulty of exploring peptide variants far from the wild-type sequences. If
any positive designs do exist at these greater distances, they represent an exponentially smaller
fraction of the sequence space, and when in complex with the TCR-MHC, are likely to adopt a
different binding mode compared to that of the wild-type.

Minor adjustments to the design protocols, aimed at better exploring the design space, may have also
contributed to the observed differences in performance across the systems. For NY-ESO, which
achieved the highest design success rate, peptides with highest PAE scores were chosen for
experimental validation, whereas in EBV and MAGE, peptides were sampled uniformly above the set PAE
thresholds. In MAGE, fixing the amino-acid at position 1 to a Glutamic Acid (E) for the one-third of
the designs led to an elevated success rate in this subset
(Fig.~\ref{fig:design_experimental_mage}D); E$_1$ is shared between the two wild-type peptides and
this design choice allowed us to assess the importance of this amino acid for function. \\

\noindent{\bf The design potential of different HERMES models.} In panel A of
Figures~\ref{fig:design_experimental_nyeso}~(NY-ESO),~\ref{fig:design_experimental_ebv}~(EBV),
and~\ref{fig:design_experimental_mage}~(MAGE), we show the fidelity of designs using the four HERMES
models: HERMES-{\em fixed} with no noise, HERMES-{\em fixed} with 0.5~\AA~noise, HERMES-{\em
relaxed} with no noise, and HERMES-{\em relaxed} with 0.5~\AA~noise. Although PyRosetta relaxations
allows HERMES-{\em relaxed} to explore the sequence space more broadly, the resulting designs that
are made at greater distances from the wild-types seldom succeed in activating T-cells. This pattern
is also observed with using models trained with 0.5~\AA~Gaussian noise. Indeed, successful designs
lie within five amino acid substitutions from the wild-types, and for these designs, the un-noised
HERMES-{\em fixed} model demonstrates the highest overall fidelity. To better explore the peptide
energy landscape and to design sequences further from the wild-types, it may be necessary to adopt a
more specialized relaxation model that accounts for the intricacies of TCR-pMHC poses, while being
consistent with the neural network's model of protein environments.\\

\noindent{\bf Comparing TCRdock and AF3 as design filters.} Using our T-cell activity measurements
as the ground truth, we sought to quantify how well the TCRdock and AF3 PAE scores can be used as
predictors of T-cell activation and filtering mechanisms for peptide design. For this purpose, we
focus primarily on the NY-ESO system, for which we validated a balanced mix of positive and negative
designs, chosen based on their TCRdock PAE scores. Here, we found that TCRdock PAE could moderately
distinguish true positive from true negative designs (AUROCs = 0.91 and 0.88,
Fig.~\ref{fig:design_roc_curves}A). Notably, using TCRdock PAE resulted in a very low false negative
rate (0.03\%) but a considerably higher false positive rate (50\%); see
Fig.~\ref{fig:design_scatterplots}A top. AF3 PAE scores performed equivalently to TCRdock PAE in
differentiating positives from negatives (AUROCs = 0.89 and 0.94,
Fig.~\ref{fig:design_roc_curves}A), though they correlated only moderately with TCRdock’s PAE values
predictions (Spearman correlation $\rho=0.5$, Fig.~\ref{fig:design_scatterplots}C). Indeed, some
peptides deemed positive by TCRdock were classified as negative by AF3, and vice versa
(Fig.~\ref{fig:design_scatterplots}C).

When restricting the analysis to peptides that were deemed as positive based on their TCRdock and
AF3 PAE scores, we did not expect either TCRdock or AF3 PAE to effectively discriminate between
peptides that did or did not elicit a T-cell response in our experimental assay (i.e., true vs.
false positives). However, we found that both TCRdock PAE and AF3 PAE scores retained significant
predictive power in the NY-ESO system (AUROCs = 0.81 and 0.77, and AUROCs = 0.80 and 0.89,
Fig.~\ref{fig:design_roc_curves}A). In the EBV and MAGE systems, the TCRdock and AF3 PAE values for
peptides that we classified positive showed lower predictive power in distinguishing true from false
positives compared to the NY-ESO system (Fig.~\ref{fig:design_roc_curves}B-C). The PAE scores of
TCRdock and AF3 differed most markedly in the MAGE system (Spearman correlation $\rho = 0.09$ for
MAGE), with TCRdock assigning relatively unfavorable PAE scores to the wild-type peptides, whereas
AF3 ranked those same peptides near the top (Fig.~\ref{fig:design_scatterplots}I).

The PAE values reported by the structure prediction algorithms have been shown to be noisy
indicators of protein folding fidelity and design quality~\cite{Roney2022-am}. While larger
experimental libraries will be needed to draw definitive conclusions, it appears that PAEs from
TCRdock and AF3 have some complementary predictive strengths, and therefore, combining their
information could lead to more robust design decisions.
Moreover, both of these scores appear to be uncorrelated or only weakly correlated with HERMES
scores across systems
(Figs.~\ref{fig:design_scatterplots},~\ref{fig:design_tcrdock_nyeso}H,~\ref{fig:design_tcrdock_ebv}H,~\ref{fig:design_tcrdock_mage}H),
suggesting that AF3 and TCRdock PAE's can serve as complementary filters to HERMES during design.\\




\noindent {\bf HERMES energy as design filter.} We further sought to assess whether HERMES energy
score itself could serve as a predictor of T-cell activation and as a filtering criterion for
design.
For each of the three systems, we scored all peptides we designed with all four combinations of
HERMES protocols, regardless of which protocol was originially used for design. For each system, we
then computed two sets of ROC curves, separately for peptides designed with HERMES-{\em fixed} and
those designed with HERMES-{\em relaxed}.
We found that scores from the HERMES-{\em relaxed} models consistently performed poorly in
predicting T-cell activation, with only one out of six cases achieving an AUROC better than random;
HERMES-{\em relaxed} 0.00 as scoring criterion on the peptides designed with HERMES-{\em fixed} in
the EBV system showing AUROC = 0.66 (Fig.~\ref{fig:design_roc_curves}B).
In contrast, HERMES-{\em fixed} 0.00 was consistently the best filtering criterion, outperforming
even TCRdock and AF3 PAE score in several cases: in the EBV system, HERMES-{\em fixed} 0.00 scores
on peptides designed by HERMES-{\em fixed} reached AUROC~$=0.74$; in the MAGE system, AUROC were
0.65 and 0.89 for and peptides designed with HERMES-{\em fixed} and HERMES-{\em relaxed},
respectively (Fig.~\ref{fig:design_roc_curves}).
These results highlight two main takeaways: first, that HERMES-{\em fixed} 0.00 is the most reliable
model for both scoring and designing peptides to elicit T-cell response; second, selecting
high-scoring peptides using this model can improve hit rates, though potentially at the cost of
reduced diversity in the peptide pool.



\section{T-cell and MHC specificity as the entropy of the associated peptide distribution}
 We characterize the specificity of a TCR-MHC complex and that of the MHC by the diversity of
 peptides that they can recognize. For TCR-MHC specificity, we starting with a TCR-pMHC structural
 template, and use HERMES-{\em fixed} model to generates a peptide position weight matrix (PWM) that
 captures the ensemble of peptides recognized by the TCR and presented by the MHC. In Figure 5, we
 analyzed 105 human and 36 mouse MHC class I TCR-pMHC structures---curated by the TCRdock template
 database \cite{Bradley2023-rl}---which included only MHC alleles represented in at least two
 structures. For each MHC allele in this set, we then characterized the distribution of peptides
 that the MHC can present (independently of the TCR) using PWMs inferred by the MHC Motif
 Atlas \cite{Tadros2023-ij}. In both cases, we use the entropy of these PWMs, $S= -\sum_{\alpha, i}
 p_i(\alpha) \log_2p_i(\alpha)$, as a proxy for degeneracy of the target (TCR-MHC or MHC), where
 $\alpha \in \{1,\dots 20\}$ denotes the amino acid types, and $i \in\{ 1, \dots, \ell\}$ are
 positions along the peptide.
 
\section{Predicting MHC-I peptide presentation preferences with HERMES.} HERMES is designed to be
broadly applicable to diverse protein-interaction contexts. As such, we attempted to use HERMES-{\em
fixed} 0.00 model to recapitulate the ensemble of peptides presented by an MHC-I allele. To do so,
we gathered 416 peptide-MHC-I crystal structures from the TCR3D database~\cite{gowthaman_tcr3d_2019}
(Table~\ref{table:pmhc_info}). For each structure, we computed a peptide position position weight
matrix (PWM) with the HERMES-{\em fixed} 0.00 design protocol, conditioning on the surrounding MHC
structure and the masked wild-type peptide as the template. The resulting PWMs were then averaged
across complexes that shared the same MHC allele and peptide length to produce an allele-specific
peptide motif (for a given peptide length).

We compared these HERMES-generated PWMs with the corresponding PWMs from the MHC Motif
Atlas~\cite{Tadros2023-ij}, which we consider as ground truth; we used relative entropy
(KL-Divergence), $D_{KL}= \sum_{\alpha, i} p^{\text{motif}}_i(\alpha) \log_2
\frac{p^{\text{motif}}_i(\alpha)}{p^{\text{HERMES}}_i(\alpha)}$ as well as the entropy difference
between the two PWMs as metrics; here, $p_i^{(\cdot)}(\alpha)$ is the probability of amino acid
$\alpha$ at position $i$ of a peptide, given the morel specified in the superscript. When structural
coverage is extensive (roughly more than 30 structures of distinct peptides with a given MHC
allele), HERMES-derived motifs align closely with those from MHC Motif Atlas, as quantified by the
low KL-divergence and comparable entropy values (Fig.~\ref{fig:pmhc_pwm_entropy_comparisons}). Under
sparse coverage, HERMES motifs begin to underestimate the peptide diversity that can be presented by
an MHC allele (measured by motif entropy), a limitation we attribute to incomplete sampling of
peptide backbone conformations by HERMES-{\em fixed} (Fig.~\ref{fig:pmhc_pwm_entropy_comparisons}C).
Overall, we see a clear pattern that the KL-Divergence increases with decreasing number of
structural templates for a given MHC allele (Fig.~\ref{fig:pmhc_pwm_entropy_comparisons}B).
Importantly, even with only a handful of templates, HERMES still captures substantially more
diversity than a na\"ive baseline PWM generated by simply aligning the template peptides themselves
(Fig.~\ref{fig:pmhc_pwm_entropy_comparisons}A). Moreover, in almost all cases, HERMES generated PWMs
are more similar (lower KL-Divergence) to that of the MHC Motif Atlas than the na\"ive baseline PWM
are (Fig.~\ref{fig:pmhc_pwm_entropy_comparisons}A).
 
These results suggest that accurate modeling of MHC presentation with HERMES-{\em fixed} design
protocol requires both sufficient structure templates and peptide diversity. It should be noted that
MHC Motif Atlas's allele‑specific motifs--derived from the large body of peptide binding data
curated for individual MHC-I alleles--are very reliable, so we continue to use them as the benchmark
for assessing motif specificity. Developing design protocols that better sample peptide backbone
conformations remains an important avenue for future work.

